\chapter{Where it stops}

\textbf{Purpose:} State what a flattened grid cannot be asked, and what this book has not
established. \textbf{Scope:}
\texttt{examples/\allowbreak{}game\_\allowbreak{}theory/\allowbreak{}2\_\allowbreak{}partition/\allowbreak{}board\_\allowbreak{}width.\allowbreak{}py}

\section{A flattened grid does not carry one period}

Reversi returns 8 because nothing repeats along a Reversi row. That is a property of the corpus and
not a property of the reader, and a reader handed a grid with structure on both axes has no way to
know which of them it answered with.

The demonstration is a two-heap subtraction game. The value of a sum of independent games is the
exclusive or of their values. That is Sprague and Grundy's theorem and is not measured here. A square
grid of two heaps under one move set is therefore generated for nothing. Such a grid repeats down the column at the
width it was laid out at, and along the row at the game's own Grundy period. Three hundred and
fifteen of them, at five widths:

\begin{center}
\begin{tabular}{@{}lr@{}}
\toprule
the flattened answer was & rows \\
\midrule
the width & 108 \\
the game's own Grundy period & 96 \\
the width off by one, which is diagonal content & 68 \\
a combination of the two & 43 \\
\bottomrule
\end{tabular}
\end{center}

The same grids shuffled return the width on 5 of 100, the chance rate. None of the four groups above
is noise. Every one of them is a real repeat in the flattened sequence, and the detector
returns whichever family scored higher.

The off-by-one group is the clearest case. Where a grid's content is a function of the sum or the
difference of its coordinates, as a checkerboard is, the shortest repeat in the flattened sequence is
at the width minus one and not at the width, because the lag that returns the same value steps one
row down and one column back. The detector reports \(W-1\) and the detector is right. The ground
truth of \(W\) was the wrong question.

The rule that comes out is short. Read the axes separately where the caller knows there are two. A
single flattened period is ambiguous where they do not, and reporting it as the width is an
assumption about the corpus dressed as a measurement.

\section{What this book has not established}

\textbf{The series covers subtraction games only.} The workbook's own division says why: soundness
belongs to the construction and cost belongs to the domain's distribution. A result obtained inside
one distribution stays there. A subtraction game's Grundy sequence is exactly periodic after a short
preperiod, over a small alphabet, with no noise. It is the easiest case a period detector can be
shown, and three hundred and eighty three rows of it establish that the
detector is right where an answer exists and is checkable. They establish nothing about a corpus
that is noisy, or only approximately periodic, or long.

\textbf{The window test is shown against one family of counterexample.} Every aperiodic word tested
is Sturmian. A sequence aperiodic in some other way might hold its answer across widening windows and
defeat the test entirely, and nothing here rules that out. Nine slopes is nine.

\textbf{Two of the sixty three Sturmian rows are unexplained.} The continued fraction mechanism
accounts for sixty one. The two it does not are the weakest rows in the table by margin, at rank 0
and rank 8 of 63, which is consistent with them being noise and is not a demonstration that they are.

\textbf{The played board is one game.} Reversi has a flip rule, and the finding that the departure
reads the rules and not the play may be a fact about flip rules. A game whose rules constrain where
a piece is placed, and flip nothing, might well leave a trace that a stronger player deepens. Nothing
here has tried one.

\textbf{The three arms are weak.} Corner beats random by 0.645 to 0.355 and that is a real ordering,
but all three are one line of heuristic. The correlation of 0.067 is measured over a strength range
that a competent Reversi program would sit entirely above, and a wider range could show a
relationship this one cannot see.

\textbf{A game is not a natural object.} A person put its regularity there deliberately. This is a
positive control, in the same class as a published cell edge and with less excuse, since nobody had
to go and measure a subtraction game. A positive control shows what a memoryless control cannot, and
nothing larger is claimed for it here.

\section{What the book is for, if one thing}

Not the three hundred and eighty three. Those were cheap by design, and a domain that hands out its
own answers is not hard to agree with.

The chapter worth carrying forward is the chapter where the instrument was confident, cleared every
floor by three to twenty nine times, produced a number with a clean mechanism behind it, and was
still not reporting what its output was called. That could not have been found with a memoryless
control, because a memoryless control has no structure to be confidently wrong about. It took an
object with real structure and a published answer saying the structure is not a period.

The test that catches it is free and it was not in the instrument. It is now.
